\chapter{Reading and writing}
\label{ch:read-write}

\section{The Read registers tab}

With a manual request, the values come back one per row, each with its address
(\ui{List}), or ten per row (\ui{Grid}), where a value's address is the row's
first address plus the column. The choice is kept. Headers stay in view while
the values scroll.

A line above the values says whether they can be written: holding registers
and coils are \val{writable}, and their values are outlined like fields; input
registers and discrete inputs are \val{read only}, and so is everything while
read-only mode is on. Click a value to select it: it gets a frame, a \ui{▸}
next to its address in the list, and the line shows its address and value, so
the address is never in doubt in the grid either. The arrow keys move the
selection.

\screenshot[6cm]{read-registers}{Holding registers read as 32-bit values.}
  {The Read registers tab after reading 20 holding registers, shown as 32-bit
   signed in CDAB order, with the address menu open.}

\begin{description}[style=nextline, leftmargin=0pt, font=\intermedium]
  \item[Show as] Shows the registers as unsigned or signed 16-bit numbers,
    hexadecimal, binary, ASCII text, 32-bit signed or unsigned integers, or
    32-bit floats, in the byte order chosen next to it. When the right type is
    not known, the Data decoder shows them all at once, 64-bit types included
    (Section~\ref{sec:data-decoder}).
  \item[Bits] Shows the 16 bits of each register. A bit of a holding register
    can be switched on its own: Modbench writes the register with only that
    bit changed.
  \item[Addresses] The address menu chooses the notation
    (Section~\ref{sec:notation}).
  \item[List, Grid] The layout of the values.
  \item[Right-click] Copy a value, an address or all values, add the address
    to the chart, or \nav{Write here…} and \nav{Write all…}.
\end{description}

\section{Writing values}
\label{sec:write}

The \tab{Write values} tab writes coils and registers.

\subsection{Without a map}

Choose the function (single or multiple coils, single or multiple
registers), the start address and the values, then \uiaccent{Write}.
\ui{Copy data to write values} fills the list with the values last read, to
change one and write them back.

\subsection{With a map}

The map's writable registers are listed by name, each with a control for its
type: a number in engineering units, a list of the map's labels for state
registers, a switch for bits. Modbench applies the inverse of the scaling,
checks the type's range and the map's limits, and only then writes.

\screenshot[6cm]{write-values}{Writable registers of a map; one mode just written.}
  {The Write values tab with a map: registers with their limits, a text field,
   and a mode chosen from its labels and written.}

\begin{description}[style=nextline, leftmargin=0pt, font=\intermedium]
  \item[Limits] A value outside the map's limits is refused before anything
    is sent, with the limit it broke.
  \item[Bit fields] Writing a bit or a field reads the register first and
    changes only its bits.
  \item[Custom values] For a register with labels, a value that is not one of
    them can still be typed.
\end{description}

\subsection{Templates}

A template copies the settings of one device to another of the same model.
\ui{Save template…} saves the values last read from the writable registers;
\ui{Load template…} writes the values of a saved template to the device shown.

\section{Writing safely}

Writes change how a real machine behaves. Modbench has three guards, in
\nav{Settings\then Connection}:

\begin{description}[style=nextline, leftmargin=0pt, font=\intermedium]
  \item[Confirm before writing to a device] Asks before every write: write
    commands, map values and templates.
  \item[Read-only mode] Blocks every write, also on a connection that is
    already open. Step 3 shows a lock while it is on; click it to go to the
    setting.
  \item[Map mismatch] After a device does not match its map, writes always ask
    first.
\end{description}

\begin{warning}
A write is sent as soon as it is confirmed. Check the unit ID: on a serial
line, unit ID \val{0} is a broadcast that every device takes.
\end{warning}
